% filtration.tex — Vietoris-Rips filtration on a small point cloud
% Include via % filtration.tex — Vietoris-Rips filtration on a small point cloud
% Include via % filtration.tex — Vietoris-Rips filtration on a small point cloud
% Include via % filtration.tex — Vietoris-Rips filtration on a small point cloud
% Include via \input{figures/filtration} from overview.tex

\begin{figure}[H]
\centering
\begin{tikzpicture}[
    point/.style={circle, fill=black, inner sep=1.8pt},
    edge/.style={thick, blue!60!black},
    tri/.style={fill=blue!15, draw=none},
    panel label/.style={font=\small\bfseries, anchor=north},
    eps label/.style={font=\small, anchor=north},
    >=Stealth
]

% --- Point coordinates (shared across all four panels) ---
% Six points arranged in a rough cluster with one outlier
\newcommand{\drawpoints}[1]{
    \node[point] (A#1) at ($(#1,0) + (0.0, 0.8)$) {};
    \node[point] (B#1) at ($(#1,0) + (0.7, 1.5)$) {};
    \node[point] (C#1) at ($(#1,0) + (1.4, 0.9)$) {};
    \node[point] (D#1) at ($(#1,0) + (0.3, 0.0)$) {};
    \node[point] (E#1) at ($(#1,0) + (1.1, 0.1)$) {};
    \node[point] (F#1) at ($(#1,0) + (2.6, 1.2)$) {};
}

% ---- Panel 1: points only, varepsilon = 0 ----
\drawpoints{0}
\node[panel label] at (0.8, -0.5) {(a)};
\node[eps label] at (0.8, -0.9) {$\varepsilon = 0$};

% ---- Panel 2: a few short edges, varepsilon = 0.9 ----
\drawpoints{4}
\draw[edge] (A4) -- (D4);
\draw[edge] (D4) -- (E4);
\draw[edge] (A4) -- (B4);
\node[panel label] at (4.8, -0.5) {(b)};
\node[eps label] at (4.8, -0.9) {$\varepsilon = 0.9$};

% ---- Panel 3: more edges + a filled triangle, varepsilon = 1.2 ----
\drawpoints{8}
% filled triangle A-D-E
\fill[tri] (A8.center) -- (D8.center) -- (E8.center) -- cycle;
\draw[edge] (A8) -- (D8);
\draw[edge] (D8) -- (E8);
\draw[edge] (A8) -- (E8);
\draw[edge] (A8) -- (B8);
\draw[edge] (A8) -- (C8);
\draw[edge] (B8) -- (C8);
\draw[edge] (C8) -- (E8);
\node[panel label] at (8.8, -0.5) {(c)};
\node[eps label] at (8.8, -0.9) {$\varepsilon = 1.2$};

% ---- Panel 4: all edges, varepsilon = 2.0 ----
\drawpoints{12}
% filled triangles
\fill[tri] (A12.center) -- (D12.center) -- (E12.center) -- cycle;
\fill[tri] (A12.center) -- (B12.center) -- (C12.center) -- cycle;
\fill[tri] (A12.center) -- (C12.center) -- (E12.center) -- cycle;
\fill[tri] (C12.center) -- (E12.center) -- (F12.center) -- cycle;
% all edges
\draw[edge] (A12) -- (B12);
\draw[edge] (A12) -- (C12);
\draw[edge] (A12) -- (D12);
\draw[edge] (A12) -- (E12);
\draw[edge] (B12) -- (C12);
\draw[edge] (B12) -- (F12);
\draw[edge] (C12) -- (E12);
\draw[edge] (C12) -- (F12);
\draw[edge] (D12) -- (E12);
\draw[edge] (E12) -- (F12);
\node[panel label] at (12.8, -0.5) {(d)};
\node[eps label] at (12.8, -0.9) {$\varepsilon = 2.0$};

% ---- Arrow showing direction ----
\draw[->, thick, gray] (2.2, -1.4) -- (13.8, -1.4)
    node[midway, below, font=\small, gray] {$\varepsilon$ increases};

\end{tikzpicture}
\caption{%
    Vietoris-Rips filtration on six points at increasing thresholds
    $\varepsilon$. At $\varepsilon = 0$ only the points exist. As
    $\varepsilon$ grows, edges appear between points whose distance is at
    most $\varepsilon$, and filled triangles (2-simplices) appear when all
    three pairwise edges are present. Tracking how connected components
    merge ($H_0$) and loops form and fill ($H_1$) across the filtration is
    the core idea of persistent homology.%
}
\label{fig:filtration}
\end{figure}
 from overview.tex

\begin{figure}[H]
\centering
\begin{tikzpicture}[
    point/.style={circle, fill=black, inner sep=1.8pt},
    edge/.style={thick, blue!60!black},
    tri/.style={fill=blue!15, draw=none},
    panel label/.style={font=\small\bfseries, anchor=north},
    eps label/.style={font=\small, anchor=north},
    >=Stealth
]

% --- Point coordinates (shared across all four panels) ---
% Six points arranged in a rough cluster with one outlier
\newcommand{\drawpoints}[1]{
    \node[point] (A#1) at ($(#1,0) + (0.0, 0.8)$) {};
    \node[point] (B#1) at ($(#1,0) + (0.7, 1.5)$) {};
    \node[point] (C#1) at ($(#1,0) + (1.4, 0.9)$) {};
    \node[point] (D#1) at ($(#1,0) + (0.3, 0.0)$) {};
    \node[point] (E#1) at ($(#1,0) + (1.1, 0.1)$) {};
    \node[point] (F#1) at ($(#1,0) + (2.6, 1.2)$) {};
}

% ---- Panel 1: points only, varepsilon = 0 ----
\drawpoints{0}
\node[panel label] at (0.8, -0.5) {(a)};
\node[eps label] at (0.8, -0.9) {$\varepsilon = 0$};

% ---- Panel 2: a few short edges, varepsilon = 0.9 ----
\drawpoints{4}
\draw[edge] (A4) -- (D4);
\draw[edge] (D4) -- (E4);
\draw[edge] (A4) -- (B4);
\node[panel label] at (4.8, -0.5) {(b)};
\node[eps label] at (4.8, -0.9) {$\varepsilon = 0.9$};

% ---- Panel 3: more edges + a filled triangle, varepsilon = 1.2 ----
\drawpoints{8}
% filled triangle A-D-E
\fill[tri] (A8.center) -- (D8.center) -- (E8.center) -- cycle;
\draw[edge] (A8) -- (D8);
\draw[edge] (D8) -- (E8);
\draw[edge] (A8) -- (E8);
\draw[edge] (A8) -- (B8);
\draw[edge] (A8) -- (C8);
\draw[edge] (B8) -- (C8);
\draw[edge] (C8) -- (E8);
\node[panel label] at (8.8, -0.5) {(c)};
\node[eps label] at (8.8, -0.9) {$\varepsilon = 1.2$};

% ---- Panel 4: all edges, varepsilon = 2.0 ----
\drawpoints{12}
% filled triangles
\fill[tri] (A12.center) -- (D12.center) -- (E12.center) -- cycle;
\fill[tri] (A12.center) -- (B12.center) -- (C12.center) -- cycle;
\fill[tri] (A12.center) -- (C12.center) -- (E12.center) -- cycle;
\fill[tri] (C12.center) -- (E12.center) -- (F12.center) -- cycle;
% all edges
\draw[edge] (A12) -- (B12);
\draw[edge] (A12) -- (C12);
\draw[edge] (A12) -- (D12);
\draw[edge] (A12) -- (E12);
\draw[edge] (B12) -- (C12);
\draw[edge] (B12) -- (F12);
\draw[edge] (C12) -- (E12);
\draw[edge] (C12) -- (F12);
\draw[edge] (D12) -- (E12);
\draw[edge] (E12) -- (F12);
\node[panel label] at (12.8, -0.5) {(d)};
\node[eps label] at (12.8, -0.9) {$\varepsilon = 2.0$};

% ---- Arrow showing direction ----
\draw[->, thick, gray] (2.2, -1.4) -- (13.8, -1.4)
    node[midway, below, font=\small, gray] {$\varepsilon$ increases};

\end{tikzpicture}
\caption{%
    Vietoris-Rips filtration on six points at increasing thresholds
    $\varepsilon$. At $\varepsilon = 0$ only the points exist. As
    $\varepsilon$ grows, edges appear between points whose distance is at
    most $\varepsilon$, and filled triangles (2-simplices) appear when all
    three pairwise edges are present. Tracking how connected components
    merge ($H_0$) and loops form and fill ($H_1$) across the filtration is
    the core idea of persistent homology.%
}
\label{fig:filtration}
\end{figure}
 from overview.tex

\begin{figure}[H]
\centering
\begin{tikzpicture}[
    point/.style={circle, fill=black, inner sep=1.8pt},
    edge/.style={thick, blue!60!black},
    tri/.style={fill=blue!15, draw=none},
    panel label/.style={font=\small\bfseries, anchor=north},
    eps label/.style={font=\small, anchor=north},
    >=Stealth
]

% --- Point coordinates (shared across all four panels) ---
% Six points arranged in a rough cluster with one outlier
\newcommand{\drawpoints}[1]{
    \node[point] (A#1) at ($(#1,0) + (0.0, 0.8)$) {};
    \node[point] (B#1) at ($(#1,0) + (0.7, 1.5)$) {};
    \node[point] (C#1) at ($(#1,0) + (1.4, 0.9)$) {};
    \node[point] (D#1) at ($(#1,0) + (0.3, 0.0)$) {};
    \node[point] (E#1) at ($(#1,0) + (1.1, 0.1)$) {};
    \node[point] (F#1) at ($(#1,0) + (2.6, 1.2)$) {};
}

% ---- Panel 1: points only, varepsilon = 0 ----
\drawpoints{0}
\node[panel label] at (0.8, -0.5) {(a)};
\node[eps label] at (0.8, -0.9) {$\varepsilon = 0$};

% ---- Panel 2: a few short edges, varepsilon = 0.9 ----
\drawpoints{4}
\draw[edge] (A4) -- (D4);
\draw[edge] (D4) -- (E4);
\draw[edge] (A4) -- (B4);
\node[panel label] at (4.8, -0.5) {(b)};
\node[eps label] at (4.8, -0.9) {$\varepsilon = 0.9$};

% ---- Panel 3: more edges + a filled triangle, varepsilon = 1.2 ----
\drawpoints{8}
% filled triangle A-D-E
\fill[tri] (A8.center) -- (D8.center) -- (E8.center) -- cycle;
\draw[edge] (A8) -- (D8);
\draw[edge] (D8) -- (E8);
\draw[edge] (A8) -- (E8);
\draw[edge] (A8) -- (B8);
\draw[edge] (A8) -- (C8);
\draw[edge] (B8) -- (C8);
\draw[edge] (C8) -- (E8);
\node[panel label] at (8.8, -0.5) {(c)};
\node[eps label] at (8.8, -0.9) {$\varepsilon = 1.2$};

% ---- Panel 4: all edges, varepsilon = 2.0 ----
\drawpoints{12}
% filled triangles
\fill[tri] (A12.center) -- (D12.center) -- (E12.center) -- cycle;
\fill[tri] (A12.center) -- (B12.center) -- (C12.center) -- cycle;
\fill[tri] (A12.center) -- (C12.center) -- (E12.center) -- cycle;
\fill[tri] (C12.center) -- (E12.center) -- (F12.center) -- cycle;
% all edges
\draw[edge] (A12) -- (B12);
\draw[edge] (A12) -- (C12);
\draw[edge] (A12) -- (D12);
\draw[edge] (A12) -- (E12);
\draw[edge] (B12) -- (C12);
\draw[edge] (B12) -- (F12);
\draw[edge] (C12) -- (E12);
\draw[edge] (C12) -- (F12);
\draw[edge] (D12) -- (E12);
\draw[edge] (E12) -- (F12);
\node[panel label] at (12.8, -0.5) {(d)};
\node[eps label] at (12.8, -0.9) {$\varepsilon = 2.0$};

% ---- Arrow showing direction ----
\draw[->, thick, gray] (2.2, -1.4) -- (13.8, -1.4)
    node[midway, below, font=\small, gray] {$\varepsilon$ increases};

\end{tikzpicture}
\caption{%
    Vietoris-Rips filtration on six points at increasing thresholds
    $\varepsilon$. At $\varepsilon = 0$ only the points exist. As
    $\varepsilon$ grows, edges appear between points whose distance is at
    most $\varepsilon$, and filled triangles (2-simplices) appear when all
    three pairwise edges are present. Tracking how connected components
    merge ($H_0$) and loops form and fill ($H_1$) across the filtration is
    the core idea of persistent homology.%
}
\label{fig:filtration}
\end{figure}
 from overview.tex

\begin{figure}[H]
\centering
\begin{tikzpicture}[
    point/.style={circle, fill=black, inner sep=1.8pt},
    edge/.style={thick, blue!60!black},
    tri/.style={fill=blue!15, draw=none},
    panel label/.style={font=\small\bfseries, anchor=north},
    eps label/.style={font=\small, anchor=north},
    >=Stealth
]

% --- Point coordinates (shared across all four panels) ---
% Six points arranged in a rough cluster with one outlier
\newcommand{\drawpoints}[1]{
    \node[point] (A#1) at ($(#1,0) + (0.0, 0.8)$) {};
    \node[point] (B#1) at ($(#1,0) + (0.7, 1.5)$) {};
    \node[point] (C#1) at ($(#1,0) + (1.4, 0.9)$) {};
    \node[point] (D#1) at ($(#1,0) + (0.3, 0.0)$) {};
    \node[point] (E#1) at ($(#1,0) + (1.1, 0.1)$) {};
    \node[point] (F#1) at ($(#1,0) + (2.6, 1.2)$) {};
}

% ---- Panel 1: points only, varepsilon = 0 ----
\drawpoints{0}
\node[panel label] at (0.8, -0.5) {(a)};
\node[eps label] at (0.8, -0.9) {$\varepsilon = 0$};

% ---- Panel 2: a few short edges, varepsilon = 0.9 ----
\drawpoints{4}
\draw[edge] (A4) -- (D4);
\draw[edge] (D4) -- (E4);
\draw[edge] (A4) -- (B4);
\node[panel label] at (4.8, -0.5) {(b)};
\node[eps label] at (4.8, -0.9) {$\varepsilon = 0.9$};

% ---- Panel 3: more edges + a filled triangle, varepsilon = 1.2 ----
\drawpoints{8}
% filled triangle A-D-E
\fill[tri] (A8.center) -- (D8.center) -- (E8.center) -- cycle;
\draw[edge] (A8) -- (D8);
\draw[edge] (D8) -- (E8);
\draw[edge] (A8) -- (E8);
\draw[edge] (A8) -- (B8);
\draw[edge] (A8) -- (C8);
\draw[edge] (B8) -- (C8);
\draw[edge] (C8) -- (E8);
\node[panel label] at (8.8, -0.5) {(c)};
\node[eps label] at (8.8, -0.9) {$\varepsilon = 1.2$};

% ---- Panel 4: all edges, varepsilon = 2.0 ----
\drawpoints{12}
% filled triangles
\fill[tri] (A12.center) -- (D12.center) -- (E12.center) -- cycle;
\fill[tri] (A12.center) -- (B12.center) -- (C12.center) -- cycle;
\fill[tri] (A12.center) -- (C12.center) -- (E12.center) -- cycle;
\fill[tri] (C12.center) -- (E12.center) -- (F12.center) -- cycle;
% all edges
\draw[edge] (A12) -- (B12);
\draw[edge] (A12) -- (C12);
\draw[edge] (A12) -- (D12);
\draw[edge] (A12) -- (E12);
\draw[edge] (B12) -- (C12);
\draw[edge] (B12) -- (F12);
\draw[edge] (C12) -- (E12);
\draw[edge] (C12) -- (F12);
\draw[edge] (D12) -- (E12);
\draw[edge] (E12) -- (F12);
\node[panel label] at (12.8, -0.5) {(d)};
\node[eps label] at (12.8, -0.9) {$\varepsilon = 2.0$};

% ---- Arrow showing direction ----
\draw[->, thick, gray] (2.2, -1.4) -- (13.8, -1.4)
    node[midway, below, font=\small, gray] {$\varepsilon$ increases};

\end{tikzpicture}
\caption{%
    Vietoris-Rips filtration on six points at increasing thresholds
    $\varepsilon$. At $\varepsilon = 0$ only the points exist. As
    $\varepsilon$ grows, edges appear between points whose distance is at
    most $\varepsilon$, and filled triangles (2-simplices) appear when all
    three pairwise edges are present. Tracking how connected components
    merge ($H_0$) and loops form and fill ($H_1$) across the filtration is
    the core idea of persistent homology.%
}
\label{fig:filtration}
\end{figure}
